%%%%%%%%%%%%%%%%%%%%%%%%%%%%%%%%%%%%%%%%%%%%%%%%%
\section{坐标空间准 PDF}
\label{sec:def}

研究准 PDF 的兴趣与热情，源于从 LQCD 第一性原理计算中提取 PDF 与强子结构的前景。因此，研究准 PDF 的价值取决于准 PDF 与 PDF 之间匹配的可靠性与精度，如文献 \cite{Ji:2013dva} 式 (11) 所提出的
\begin{equation}
\label{eq:matching_ji}
\tilde{q}(x,\mu^2,P^z) = \int_x^1 \frac{dy}{y} Z\left(\frac{x}{y},\frac{\mu}{P^z}\right) q(y,\mu^2) \, ,
\end{equation}
对夸克分布，其中 $Z(x,\mu/P^z)=\delta(x-1)+(\alpha_s/2\pi)Z^{(1)}(x,\mu/P^z) +\dots \to \delta(x-1)$（当 $P^z\to \infty$）。这等价于要求准 PDF 的 PQCD 因子化具有可靠性与精度：
\begin{eqnarray}
\label{eq:factorization}
\tilde{f}_{i/p}(\tilde{x},\tilde{\mu}^2,p_z) & \approx &
\sum_{j} \int_0^1 \frac{dx}{x} {\cal C}_{ij}(\frac{\tilde{x}}{x},\tilde{\mu}^2,\mu^2,p_z)
\nonumber \\
&& \hskip 0.3in
\times f_{j/p}(x,\mu^2) + {\cal O}\left(\frac{1}{\tilde{\mu}^2}\right),
\end{eqnarray}
其中 $i, j =q, \bar{q}, g$，${\cal C}$ 是红外安全且微扰可计算的匹配系数，$\tilde{\mu}$ 是准 PDF 的重正化标度，$\mu (\sim{\hskip -0.05in}\tilde{\mu})$ 是 PDF 的因子化标度，而幂次修正 ${\cal O}(1/\tilde{\mu}^2)$ 由高扭度夸克--胶子关联函数的大小刻画，一如所有共线 PQCD 因子化框架。我们发现，若把式~(\ref{eq:factorization}) 左边准 PDF 的微扰紫外发散用动量截断正规化，则右边的因子化卷积表现良好。然而，动量截断正规子在微扰展开的高阶计算中难以一致地实施。另一方面，我们注意到，若用维数正规化（DR）正规化准 PDF 的微扰紫外发散，则式~\eqref{eq:factorization} 右边对动量分数 $x$ 的积分会发散。这是因为海夸克与胶子 PDF 当 $x\to0$ 时满足 $f_{j/p}(x,\mu^2) \to x^{-\alpha}$（$1<\alpha<2$），而用 DR 计算出的匹配系数 ${\cal C}^{(1)}_{ij}(\tilde{x}/x,\tilde{\mu}^2,\mu^2,p_z)$ 在 $x\to 0$ 时含有正比于 $1/(\tilde{x}/x)$ 的项 \cite{Ma:2014jla,Xiong:2013bka}。结果，式~(\ref{eq:factorization}) 中 $x$ 积分的下端导致发散 $\int_0 dx/x\, (x/\tilde{x})\, x^{-\alpha} \to \infty$。

把式~(\ref{eq:factorization}) 的因子化公式应用于一个渐近部分子态``$j$''，并将因子化等式两边按 $\alpha_s$ 幂次展开到一阶，得到 ${\cal C}^{(1)}_{ij}(y,\tilde{\mu}^2,\mu^2,p_z)=\tilde{f}^{(1)}_{i/j}(y,\tilde{\mu}^2,p_z) - f^{(1)}_{i/j}(y,\mu^2)$。由于 $f^{(1)}_{i/j}(y,\mu^2)$ 在 $y\to\infty$ 时消失，${\cal C}^{(1)}_{ij}(\tilde{x}/x,\tilde{\mu}^2,\mu^2,p_z)$ 在 $x\to 0$ 的渐近行为完全由 $\tilde{f}^{(1)}_{i/j}(\tilde{x},\tilde{\mu}^2,p_z)$ 的大 $\tilde{x}$ 行为决定。

当 $\tilde{x}\to\infty$ 时，由式~\eqref{eq:ftq} 与 \eqref{eq:ftg} 定义的准 PDF 的大动量行为，与两场间隔 $\xi_z\to 0$ 时强子矩阵元的渐近行为密切相关。上文发现的式~(\ref{eq:factorization}) 因子化形式中 $x$ 卷积的发散表明：当 $\xi_z\to 0$ 时强子矩阵元可能不适定——下一节将证明确实如此。正如文献 \cite{Ma:2014jla} 所指出的，坐标空间准 PDF 对 LQCD 计算、对重正化讨论以及对 PDF 提取都是更好的量。

我们如下定义坐标空间准 PDF：
\begin{align}\label{eq:ftqx}
\begin{split}
&\tilde{F}_{q/p}(\xi_z,\tilde{\mu}^2,p_z)\\
=& \frac{e^{i p_z \xi_z}}{p_z} \langle h(p)| \overline{\psi}_q(\xi_z)\,\frac{\gamma_z}{2}
\Phi_{n_z}^{(f)}(\{\xi_z,0\})\, \psi_q(0) | h(p) \rangle,
\end{split}
\end{align}
以及
\begin{align}\label{eq:ftgx}
\begin{split}
&\tilde{F}_{g/p}(\xi_z,\tilde{\mu}^2,p_z)\\
=&\frac{e^{i p_z \xi_z}}{p_z^2}\langle h(p)| F_{z}^{\ \nu}(\xi_z)\,
\Phi_{n_z}^{(a)}(\{\xi_z,0\})\, F_{z\nu}(0) | h(p) \rangle,
\end{split}
\end{align}
其中 $\tilde{\mu}$ 是重正化标度，$\Phi_{n_z}^{(f,a)}(\{\xi_z,0\})$ 是第 \ref{sec:intro} 节定义的规范链接。我们还定义 $n_z^\mu = (0,0_\perp, 1)$，度规 $g^{\mu\nu}=\text{diag}(1,-1,-1,-1)$，并对任意矢量 $v^\mu$ 有 $v\cdot n_z =- v_z$，于是 $n_z^2 = -1$。


我们将在下一节证明：坐标空间准 PDF 对有限的 $\xi_z$ 是良定的，而当 $\xi_z\to0$ 时发散。如果想要良定的动量空间准 PDF，就必须恰当而一致地扣除 $\xi_z\to0$ 处的发散。由准 PDF 的定义可得
\begin{align}
\tilde{F}_{i/p}(-\xi_z,\tilde{\mu}^2,p_z)=\tilde{F}_{i/p}^*(\xi_z,\tilde{\mu}^2,p_z).
\end{align}
为了利用这一关系并简化讨论，下面的计算假设 $\xi_z>0$。对于最终结果，我们会把它写成对任意 $\xi_z$ 取值都成立的形式。

在进入证明准 PDF 可重正性的细节之前，先简要说明准 PDF 重正化的复杂性：\\
{\bf 1) 洛伦兹对称性破缺。} 由于准 PDF 明显依赖``$z$''方向，为辨认所有可能的紫外发散，需要对每个圈分别研究四维圈动量积分与三维圈动量积分（固定``$z$''方向）\footnote{注意，只有这两种情形需要研究；其他圈动量积分情形不会产生新的紫外发散。}。对一个 $n$ 圈费曼图而言，这意味着 $2^n$ 种不同的情形，难以处理。\\
{\bf 2) 需要复合算符的重正化。} 举例来说，选取 $A_z=0$ 的轴向规范，夸克准 PDF 变为
\begin{align}\label{eq:ftqxAz}
\begin{split}
&\tilde{F}_{q/p}(\xi_z,\tilde{\mu}^2,p_z)= \frac{e^{i p_z \xi_z}}{p_z} \left. \langle h(p)| \overline{\psi}_q(\xi_z)\,\frac{\gamma_z}{2}\, \psi_q(0) | h(p) \rangle\right. .
\end{split}
\end{align}
要重正化夸克准 PDF，既要重正化式~\eqref{eq:ftqxAz} 中的单个夸克场 $\overline{\psi}_q(\xi_z)$ 与 $\psi_q(0)$，也要在双定域复合算符整体产生紫外发散时将其作为整体重正化。场的重正化可以通过使用 $A_z=0$ 规范下重正化的 QCD 拉氏量自然处理；而双定域复合算符的重正化不在 QCD 重正化的覆盖范围之内。这与证明 PDF 可重正性的程序实质相同：例如对 $A^+=0$ 规范中的夸克 PDF，必须验证定义夸克 PDF 的双定域复合算符的乘法重正化，才能证明其可重正性。正是双定域复合算符的重正化使夸克 PDF 与胶子 PDF 混合 \cite{Collins:1981uw}。也就是说，仅凭 QCD 拉氏量本身的可重正性，不足以保证由复合算符定义的准 PDF 的可重正性。


在本文余下部分的显式计算与讨论中，我们选取费曼规范，因为费曼规范下 QCD 拉氏量的重正化是众所周知的。既然单个夸克场与胶子场的重正化是良定的，我们将聚焦于定义准 PDF 的复合算符的重正化。下面主要讨论坐标空间准 PDF，因此凡提到准 PDF，均指坐标空间准 PDF。我们先集中讨论夸克准 PDF，最后再把研究推广到胶子准 PDF。
